\section{Random variables and their distributions}

As we have seen, a probability space is an abstract
concept that represents our intuitive notion of a
probabilistic experiment. Such an experiment however
can be very long (even infinite) and contain a lot of
information. To make things more manageable, we
consider numerical-valued functions on probability
spaces, which we call \textbf{random variables}.
However, not any function will do: a random variable
has to relate in a nice way to the measurable space
structure, so that we will be able to ask questions
like ``what is the probability that this random
variable takes a value less than 8'', etc. This leads
us to the following definitions.

\begin{defn}
If $(\Omega_1, {\cal F}_1)$ and
$(\Omega_2, {\cal F}_2)$ are two measurable spaces, a
function $X:\Omega_1\to \Omega_2$ is called
$\textbf{measurable}$ if for any set
$E \in {\cal F}_2$, the set

\[
X^{-1}(E) = \{ \omega \in \Omega_1 : X(\omega) \in E \}
\]

\noindent is in ${\cal F}_1$.
\end{defn}

\begin{defn}
If $(\Omega, {\cal F}, \mathrm{P})$ is a probability
space, a real-valued function $X:\Omega \to \R$ is
called a $\textbf{random variable}$ if it is a
measurable function when considered as a function from
the measurable space $(\Omega, {\cal F})$ to the
measurable space $(\R, {\cal B})$, where ${\cal B}$
is the Borel $\sigma$-algebra on $\R$, namely the
$\sigma$-algebra generated by the intervals.
\end{defn}

\begin{xexercise}
\label{ex-checkrv}
Let $(\Omega_1, {\cal F}_1)$ and
$(\Omega_2, {\cal F}_2)$ be two measurable spaces
such that ${\cal F}_2$ is the $\sigma$-algebra
generated by a collection ${\cal A}$ of subsets of
$\Omega_2$. Prove that a function
$X:\Omega_1 \to \Omega_2$ is measurable if and only
if $X^{-1}(A) \in {\cal F}_1$ for all
$A \in {\cal A}$.
\end{xexercise}

\noindent It follows that the random variables are exactly
those real-valued functions on $\Omega$ for which the
question

\[
\textrm{``What is the probability that $a < X < b$?''}
\]

\noindent has a well-defined answer for all $a<b$. This
observation makes it easier in practice to check if a
given function is a random variable or not, since
working with intervals is much easier than with the
rather unwieldy (and mysterious, until you get used to
them) Borel sets.

What can we say about the behavior of a random
variable $X$ defined on a probability space
$(\Omega, {\cal F}, \mathrm{P})$? All the information
is contained in a new probability measure $\mu_X$ on
the measurable space $(\R, {\cal B})$ that is induced
by $X$, defined by

\[
\mu_X(A) = \mathrm{P}(X^{-1}(A)) = \mathrm{P}(\omega \in \Omega : X(\omega) \in A).
\]

\noindent The number $\mu_X(A)$ is the probability that $X$
``falls in $A$'' (or ``takes its value in $A$'').

\begin{xexercise}
Verify that $\mu_X$ is a probability measure on
$(\R, {\cal B})$. This measure is called the
\textbf{distribution} of $X$, or sometimes referred
to more fancifully as the $\textbf{law}$ of $X$. In
some textbooks it is denoted ${\cal L}_X$.
\end{xexercise}

\begin{defn}
If $X$ and $Y$ are two random variables (possibly
defined on different probability spaces), we say that
$X$ and $Y$ are $\textbf{identically distributed}$
(or \textbf{equal in distribution}) if $\mu_X = \mu_Y$
(meaning that $\mu_X(A) = \mu_Y(A)$ for any Borel set
$A\subset \R$). We denote this

\[
X \stackrel{d}{=} Y.
\]

\end{defn}

\noindent How can we check if two random variables are
identically distributed? Once again, working with
Borel sets can be difficult, but since the Borel sets
are generated by the intervals, a simpler criterion
involving just this generating family of sets exists.
The following lemma is a consequence of basic facts in
measure theory, which can be found in the Measure
Theory appendix in [Dur2010].

\begin{lem}
Two probability measures $\mu_1, \mu_2$ on the
measurable space $(\R, {\cal B})$ are equal if only
if they are equal on the generating set of intervals,
namely if

\[
\mu_1\big((a,b)\big) = \mu_2\big((a,b)\big)
\]

\noindent for all $-\infty < a<b < \infty$.
\end{lem}
